\chapter{Fondements techniques et état de l'art}
\label{chap:etat-art}

\section*{Introduction}

Ce chapitre expose les briques techniques mobilisées par le projet. Il part du problème central
--- pourquoi un grand modèle de langage seul ne peut pas répondre de façon fiable à une question
juridique --- puis présente l'architecture RAG comme réponse à ce problème, avant de détailler
les deux familles de méthodes de recherche utilisées et leur combinaison.

\section{Grands modèles de langage et hallucination}

Un grand modèle de langage (\textit{Large Language Model}, LLM) est un modèle statistique entraîné
à prédire la suite la plus plausible d'une séquence de texte. Cette formulation explique à la fois
sa fluidité et sa principale faiblesse : le modèle produit ce qui est \emph{vraisemblable}, non ce
qui est \emph{vrai}.

Ce comportement porte le nom d'\textbf{hallucination}. Dans un contexte juridique, il est
particulièrement dangereux : un modèle sollicité sur un article de loi qu'il ne connaît pas
produira volontiers un numéro d'article plausible accompagné d'un contenu inventé. La réponse
paraît crédible, la référence semble vérifiable, et rien dans la forme ne signale l'invention.

Deux stratégies existent pour ancrer un modèle dans des faits réels :

\begin{itemize}
    \item Le \textbf{\textit{fine-tuning}}, qui consiste à réentraîner le modèle sur un corpus
    spécialisé. Coûteux, difficile à mettre à jour, et surtout : le modèle continue de générer
    depuis ses poids, sans possibilité de citer une source vérifiable.
    \item La \textbf{génération augmentée par la recherche} (RAG), qui fournit au modèle les
    documents pertinents \emph{au moment de la question} et lui demande de s'appuyer uniquement
    sur eux.
\end{itemize}

C'est la seconde qui a été retenue : elle permet la citation, donc la vérification, ce qui est la
condition même de la fiabilité recherchée.

\section{L'architecture RAG}

Le principe du RAG consiste à séparer deux préoccupations habituellement confondues : \emph{savoir}
et \emph{formuler}. Le savoir est stocké dans un corpus documentaire externe et interrogeable ; le
modèle de langage n'est plus qu'un moteur de reformulation appliqué à des documents qu'on lui
fournit.

Le système se décompose en deux pipelines distincts :

\begin{itemize}
    \item Un pipeline d'\textbf{indexation}, exécuté une seule fois hors ligne : le corpus est
    découpé, vectorisé et stocké dans des index.
    \item Un pipeline de \textbf{requête}, exécuté à chaque question : la question est vectorisée,
    les passages pertinents sont retrouvés, puis transmis au modèle qui rédige la réponse.
\end{itemize}

Cette séparation a une conséquence pratique décisive pour le débogage : il est possible d'inspecter
les documents retrouvés \emph{avant} l'appel au modèle. Or la majorité des mauvaises réponses en RAG
proviennent d'une mauvaise \emph{recherche}, non d'une mauvaise \emph{génération}.

\section{Recherche sémantique par embeddings}

Un \textbf{embedding} est la représentation d'un texte sous forme de vecteur numérique de dimension
fixe, construite de telle sorte que deux textes de sens proche produisent des vecteurs proches dans
l'espace. La proximité est mesurée par la \textbf{similarité cosinus} :

\begin{equation}
    \mathrm{sim}(\vec{a}, \vec{b}) = \frac{\vec{a} \cdot \vec{b}}{\|\vec{a}\|\,\|\vec{b}\|}
\end{equation}

Pour des vecteurs normalisés, cette expression se réduit à un simple produit scalaire, borné entre
$-1$ et $1$.

L'intérêt de cette approche est qu'elle capture le \emph{sens} et non les mots : une question posée
en langage courant (« mon patron peut me virer sans raison ? ») peut retrouver un article rédigé en
langage juridique formel, sans partager aucun terme avec lui.

Sa limite est symétrique : elle est moins précise lorsque le terme exact compte --- un numéro
d'article, une expression figée comme « Loi 65-99 ».

Une \textbf{base vectorielle} stocke ces vecteurs et permet de retrouver efficacement les plus
proches voisins d'un vecteur requête.

\section{Recherche lexicale : BM25}

BM25 (\textit{Best Matching 25}) est une fonction de classement lexicale qui score un document en
fonction des termes qu'il partage avec la requête. Elle pondère chaque terme par sa rareté dans le
corpus (un mot rare est plus discriminant qu'un mot fréquent) et normalise par la longueur du
document, afin qu'un document long ne soit pas favorisé mécaniquement.

BM25 est exactement complémentaire des embeddings : excellent sur les correspondances exactes,
aveugle à la reformulation.

\section{Fusion des classements : Reciprocal Rank Fusion}

Combiner les deux méthodes pose un problème pratique : leurs scores ne sont pas comparables. Un
score BM25 n'est pas borné et peut valoir 8 ou 20 ; une similarité cosinus est bornée entre $-1$ et
$1$. Faire la moyenne de $8$ et de $0{,}7$ n'a aucun sens.

La \textbf{fusion par rang réciproque} (\textit{Reciprocal Rank Fusion}, RRF) contourne
entièrement ce problème en n'utilisant que le \emph{rang} de chaque document dans chaque liste, et
non son score brut :

\begin{equation}
    \mathrm{RRF}(d) = \sum_{m \in M} \frac{1}{k + \mathrm{rang}_m(d)}
\end{equation}

où $M$ est l'ensemble des méthodes de recherche, $\mathrm{rang}_m(d)$ le rang du document $d$ selon
la méthode $m$, et $k$ une constante d'amortissement (valeur usuelle : 60) qui limite l'influence
excessive des tout premiers rangs.

Un document bien classé par les deux méthodes obtient ainsi un score élevé, sans qu'aucune
normalisation d'échelle ne soit nécessaire.

\section{Souveraineté : la contrainte d'exécution locale}

La contrainte de souveraineté a une conséquence technique directe : \textbf{aucun composant du
cœur du système ne peut être un service en ligne}. Cela exclut d'emblée les embeddings par API,
les services de reclassement commerciaux, et les modèles de génération propriétaires accessibles
uniquement à distance.

Toute la chaîne doit donc reposer sur des modèles open-source exécutables sur du matériel
ordinaire. Cette contrainte, loin d'être seulement une limitation, structure l'ensemble des choix
techniques présentés au chapitre suivant.

\newpage
\section*{Conclusion}

Le RAG apporte une réponse directe au problème de l'invention de références : en fournissant au
modèle les textes au moment de la question, il rend la citation possible et donc la vérification
réalisable. La recherche sémantique et la recherche lexicale étant complémentaires, leur fusion
par rang réciproque permet de bénéficier des deux sans arbitrage d'échelle. Le chapitre suivant
traduit ces principes en une architecture concrète.
